\chapter{PENDAHULUAN}

\section{Latar Belakang}
Pada aplikasi yang terdiri atas beberapa layanan, satu permintaan dapat memicu pekerjaan pada proses yang berbeda. Sebagai contoh, alur pemesanan dapat menjalankan pembayaran dan reservasi persediaan secara konkuren, kemudian memulai pengiriman setelah kedua pekerjaan tersebut selesai. Penggunaan interaksi asinkron memungkinkan pekerjaan tidak harus menunggu satu rangkaian eksekusi sekuensial. Akan tetapi, pengembang perlu memastikan bahwa dependensi bisnis tetap dipenuhi meskipun waktu penyelesaian, pengiriman pesan, dan pencatatan observasi berbeda antarlayanan.

Masalah tersebut tidak dapat dijawab hanya dengan menyusun seluruh peristiwa berdasarkan waktu fisik. \citet{lamport1978} menjelaskan relasi \ist{happened-before} yang membentuk urutan parsial (\ist{partial order}) pada sistem terdistribusi. Relasi ini membedakan peristiwa yang berhubungan melalui urutan lokal atau pertukaran pesan dari peristiwa yang tidak memiliki dependensi kausal. Timestamp adalah catatan waktu fisik suatu peristiwa. Perbedaan timestamp bukan bukti yang cukup untuk menetapkan hubungan sebab-akibat lintas proses. Pada contoh pemesanan, pembayaran dan reservasi persediaan boleh selesai dalam urutan mana pun; yang perlu dipertahankan adalah dependensi keduanya terhadap pengiriman.

OpenTelemetry menyediakan model observasi berupa \ist{trace}, \ist{span}, peristiwa pada \ist{span}, dan atribut \citep{oteltraces}. Model tersebut memungkinkan data dari beberapa layanan dikaitkan, tetapi observasi teknis tidak langsung sama dengan peristiwa bisnis. \ist{Span} permintaan pembayaran, misalnya, dapat dimulai sebelum pembayaran berhasil. Karena itu, pemeriksa yang menggunakan awal \ist{span} sebagai bukti keberhasilan pembayaran berpotensi menilai alur secara keliru. Pemilihan titik peristiwa dan cara pemetaan observasi menjadi dependensi perlu dinyatakan dalam kontrak instrumentasi.

Penelitian terdahulu telah menghubungkan perilaku yang diamati dengan perilaku yang diharapkan. \citet{sakai2022} mengkaji pembentukan model proses layanan untuk pemeriksaan kesesuaian menggunakan \ist{distributed tracing}. \citet{sun2024} mengusulkan pendekatan berbantuan anotasi untuk mendeteksi inkonsistensi perilaku sistem berbasis \ist{microservices}. Pada ranah protokol terdistribusi, Ellsberg memeriksa implementasi terhadap simulasi protokol \citep{ding2025}, sedangkan TraceLink menghubungkan \ist{trace} implementasi dengan desain terverifikasi melalui ekosistem PGo \citep{hackett2025}. Kajian automata untuk kumpulan \ist{trace} lokal juga menunjukkan bahwa pemeriksaan interaksi tidak harus mengandalkan jam global bersama \citep{weilkennedy2025}.

Proposal ini berfokus pada kesesuaian alur kerja (\ist{workflow}) aplikasi terhadap mesin keadaan (\ist{state machine}) kecil yang ditulis pengembang. Telemetri dihasilkan saat aplikasi berjalan, sedangkan verifikasi dilakukan secara \ist{offline} setelah pengumpulan; penelitian ini merupakan verifikasi \ist{runtime} berbasis observasi eksekusi. Contohnya, pembayaran dan reservasi boleh muncul dalam urutan mana pun, tetapi pengiriman hanya diterima setelah keduanya dan hanya satu kali. Pemeriksa mencari sedikitnya satu urutan peristiwa yang diterima model sekaligus mempertahankan semua hubungan kausal yang diamati. Hasil kompatibel dengan model berarti terdapat urutan yang diizinkan dan sesuai bukti. Hasil ini tidak membuktikan bahwa aplikasi benar-benar menegakkan prasyarat. Jika pembayaran dan pengiriman tidak terurut dalam observasi, pemeriksa dapat menerima urutan pembayaran lebih dahulu meskipun fakta aplikasi menunjukkan pengiriman dilakukan terlalu dini; skenario S7 pada Bab III menguji keterbatasan ini.

Pemanfaatan telemetri juga memerlukan penanganan bukti yang tidak lengkap. Tidak ditemukannya pembayaran dalam suatu ekspor dapat berarti pembayaran memang tidak terjadi, tetapi dapat pula berarti peristiwa tersebut tidak berhasil dicatat atau diekspor. Mekanisme \ist{sampling} dan batas perekaman merupakan bagian yang perlu diperhatikan pada instrumentasi OpenTelemetry \citep{otelsdk}. Menyamakan kehilangan observasi dengan kegagalan aplikasi dapat menimbulkan alarm yang tidak berdasar. Oleh sebab itu, penelitian ini membedakan status pelanggaran yang terbukti dari status tidak dapat disimpulkan.

Kesenjangan yang ditelaah adalah kebutuhan akan pendekatan dengan kontrak observasi yang jelas, aturan interaksi sederhana, dan evaluasi yang memisahkan dampak urutan \ist{timestamp} dari dampak kualitas bukti. Penggunaan urutan parsial maupun pemeriksaan kesesuaian telah dikaji sebelumnya. Kontribusi yang diusulkan terletak pada perancangan dan evaluasi pemeriksaan kesesuaian alur kerja pada data OpenTelemetry untuk ruang lingkup aplikasi asinkron, termasuk batas interpretasi hubungan \ist{span} dan penjelasan peristiwa relevan ketika terjadi pelanggaran.

\section{Rumusan Masalah}
Berdasarkan latar belakang tersebut, masalah penelitian dirumuskan sebagai berikut.
\begin{enumerate}
  \item Bagaimana membentuk representasi urutan parsial dari observasi OpenTelemetry berdasarkan kontrak observasi?
  \item Bagaimana memeriksa kompatibilitas observasi dengan spesifikasi alur kerja ketika bukti tidak selalu memadai?
  \item Bagaimana pemeriksaan urutan parsial dibandingkan dengan pengurutan timestamp dalam ketepatan, cakupan, biaya, dan deteksi kesalahan aplikasi?
\end{enumerate}

\section{Batasan Masalah}
Ruang lingkup penelitian dibatasi sebagai berikut.
\begin{enumerate}
  \item Objek uji mencakup aplikasi terkontrol, alur kerja dari aplikasi \ist{microservices} sumber terbuka yang dikembangkan pihak lain, dan graf sintetis. Aplikasi eksternal dipilih berdasarkan komunikasi asinkron, kesesuaian model alur kerja, kemampuan instrumentasi, dan sumber daya yang tersedia; satu aplikasi eksternal menjadi cakupan wajib, dengan OpenTelemetry Demo rilis 2.0.2 sebagai target awal. Alur checkout--Kafka--accounting telah menjalani uji awal pengumpulan observasi. Kelayakan kasus kesalahan bisnis tetap diperiksa sebelum evaluasi resmi; revisi dan aturan penggantian dijelaskan pada Bab III.
  \item Spesifikasi utama berupa mesin keadaan berhingga untuk alur kerja terbatas per alur kerja. Penerimaan memakai semantik eksistensial atas urutan topologis sebagaimana dijelaskan pada Bab II. Penelitian tidak mengembangkan bahasa formal baru, tidak memverifikasi konsensus, dan tidak membuktikan \ist{liveness} tanpa batas.
  \item Verifikasi dilakukan secara \ist{offline} terhadap hasil observasi eksekusi. Istilah \ist{runtime} merujuk pada sumber bukti yang diperoleh saat aplikasi berjalan; pemeriksaan \ist{streaming} dan intervensi otomatis tidak termasuk sasaran awal.
  \item Masukan awal berupa ekspor OTLP JSON yang dinormalisasi. Dukungan terhadap format biner, semua varian backend tracing, dan instrumentasi otomatis seluruh aplikasi berada di luar lingkup.
  \item Identitas peristiwa, identitas alur kerja, dan semantik arah dependensi ditentukan dalam kontrak instrumentasi. \ist{Trace} standar tanpa informasi bisnis yang memadai tidak dianggap cukup untuk seluruh aturan.
  \item Tidak semua \ist{span link} diperlakukan sebagai \ist{happened-before}. Hubungan dipakai hanya jika makna kausal dan titik peristiwanya dapat dipertanggungjawabkan.
  \item Pembandingan dilakukan pada masukan dan aturan bisnis yang sama. \ist{Timestamp} digunakan hanya oleh pembanding untuk menetapkan urutan; waktu monotonic digunakan untuk pengukuran durasi verifikasi.
  \item Penjelasan pelanggaran berupa irisan peristiwa relevan yang deterministik. Minimalitas global contoh penyangkal (\ist{counterexample}) dan efektivitas penjelasan melalui studi pengguna berada di luar ruang lingkup.
\end{enumerate}

\section{Tujuan Penelitian}
Penelitian bertujuan mengembangkan dan mengevaluasi prototipe verifikasi interaksi layanan asinkron berbasis urutan parsial. Tujuan tersebut dijabarkan menjadi:
\begin{enumerate}
  \item merumuskan kontrak normalisasi observasi OpenTelemetry menjadi peristiwa dan dependensi kausal;
  \item menerapkan pencarian kesesuaian alur kerja dengan keluaran kompatibel dengan model, tidak kompatibel dengan model, dan tidak dapat disimpulkan;
  \item membandingkan ketepatan, cakupan keputusan, biaya verifikasi, dan ukuran penjelasan melalui eksperimen berlabel, serta mengevaluasi deteksi kesalahan aplikasi secara terpisah;
  \item menghasilkan penjelasan yang menampilkan aturan, alur kerja, serta peristiwa relevan bagi pelanggaran.
\end{enumerate}

\section{Manfaat Penelitian}
Secara akademis, penelitian diharapkan memberikan bukti empiris mengenai penggunaan informasi kausal dalam pemeriksaan interaksi aplikasi asinkron. Dataset berlabel dan prosedur evaluasi dapat menjadi dasar untuk menguji kembali batasan pendekatan, terutama pada hubungan antara konkurensi, perbedaan jam, dan ketidaklengkapan observasi.

Secara metodologis, kontrak observasi membantu memisahkan fakta bisnis, metadata instrumentasi, serta keputusan pemeriksa. Pemisahan tersebut memungkinkan kesalahan aplikasi dibedakan dari kesalahan ekspor atau penafsiran telemetri. Secara praktis, prototipe diharapkan membantu pengembang menelusuri ketidaksesuaian urutan pada alur kerja melalui penjelasan yang lebih terarah daripada membaca seluruh \ist{trace} secara manual.

\section{Sistematika Penulisan}
Proposal ini terdiri atas empat bab. Bab I menjelaskan latar belakang, masalah, batasan, tujuan, dan manfaat penelitian. Bab II membahas konsep yang digunakan, penelitian terkait, dan posisi penelitian. Bab III menjelaskan metode, kontrak observasi, rancangan pemeriksa, serta evaluasi dan ancaman validitas. Bab IV memuat rencana jadwal pelaksanaan penelitian. Daftar pustaka disajikan setelah Bab IV.
